\section{Evola: The Second Rumi?}

\begin{quotex}
Some souls learn nothing except from human masters; others have learned everything from invisible guides known only to themselves.
\flright{\textsc{Shihaboddin Yahya Sohravardi}}
\end{quotex}


\textbf{Fr. Frank Gelli}, a Roman by birth, now an Anglican priest living in a predominantly Muslim section of London, recently made available on Amazon an e-book, Julius Evola: The Sufi of Rome. This book tries to work on two levels. Ostensibly, it is a record of the author's many encounters with \textbf{Julius Evola}. In the 1960s, the young Fr. Gelli was a member of a rightwing movement loosely affiliated with Julius Evola. The members, including Gelli, had occasional audiences with Evola. From there, Gelli was apparently singled out for special attention and became a regular visitor to Evola's flat.

He was strangely asked to keep their relationship private, presumably not to make the others envious. Not very convincing. But why after 40 years, has Fr. Gelli broken his silence to reveal what had been sealed? We'll take a guess at that in the end.

For the second level, Gelli seems convinced that Evola is a closeted Sufi, but of a special type. A \emph{malamatiya} appears outwardly to bring blame on himself, but inwardly he is quite blameless and pious. Here Gelli is less convincing. Gelli assumes Guenon initiated Evola, but the two never met in person. He points out that Evola is rejected because of his associations with racism, anti-Semitism, and fascism. However, at the time Evola was writing on those topics, he didn't do it to bring blame on himself, rather he was trying to curry favour. Only after his side lost the war, did his positions on those issues put him on the wrong side of polite society.

It is true, nevertheless, that Evola attracts enemies and even worse friends for whatever reasons. I myself was publicly upbraided by a Sicilian professoressa from a local university when she found out I did translations of Evola. I found out later her specialty is Michel Foucault, who is a much more loathsome man. Albert Schweitzer was more explicitly racist, but is regarded as a great humanitarian; T S Eliot was a cruder anti-Semite, but that never affected his career; Martin Heidegger was publicly part of the Nazi regime, but is still considered the best philosopher of last century. The reason why Evola is such a polarizing figure is not so easy to discern.

\paragraph{The Way of Blame}


Although he claims to be praising Evola, his portrayal is so pedestrian, it has the totally opposite effect. And it cannot be due to that Sufi path. There is a story of a Zen monk who was accused by a single woman of being the father of her bastard child. He said, ``I see.'' Then he took the young boy into his home and raised him. A few years later, the real father returned and the woman at last revealed the truth. When the real father came to claim his son, the monk handed him over.

Compare that to Evola's behavior in a similar situation. Evola showed Gelli a letter from an Argentine woman who claimed she gave birth to his son. Evola protested to Gelli, doubting the proof of his paternity (that's what happens with promiscuity), and wishing she had ``done away with him''. Compare this story to the monk; unlike the monk, Evola actually was not blameless, and his character barely rises above ghetto life in the USA; certainly not what we would expect from a well-bred man.

\paragraph{Right Wing}


We learn little of the ``right wing'' group Gelli belonged to, other than they would meet in a specific café to chat, occasionally formulating a plan (usually instigated by a police informant). They were mostly thuggish types, and Gelli tells us nothing of their goals and platforms. He does reveal that it was riddled with homosexuals. Since all the loose women were on the left, Gelli often attended their rallies. He even became a Maoist in 1969, so it seems several years of conversation with Evola did not leave a lasting impression on him. Evola himself never seemed to what the right wing should be like. He did suggest that perhaps the remnants of European aristocracy may work something out.

\paragraph{Tradition}


It is likewise difficult to discern what was particularly ``Traditional'' in their chats. Evola participated in neither rites nor prayers. He never mentioned ``God'', preferring, instead, allusions to a rather vague ``transcendence''. But there was no discussion of initiation, how to reach higher states of consciousness, and so on.

Curiously, Evola told Gelli that the Roman Senate could not proceed in any course of action, until the Augurs approved it, based on their readings. But this illustrates the superiority of the spiritual authority, or sacerdotal caste, \emph{pace} his own point of view. I suppose readers can take this as Sufi wisdom; I take it as an intellectual contradiction.

\paragraph{Omar Amin von Leers}


Evola asked Gelli, on his trip to Egypt, to visit Omar Amin von Leers\footnote{\url{http://www.andrewbostom.org/blog/2011/05/30/egyptian-islamo-nazism-and-\%E2\%80\%9Comar-amin\%E2\%80\%9D-von-leers/}}, a Nazi refugee and convert to Islam. Omar mysteriously claimed Evola was ``one of us''. Gelli assumes that means Evola was a Muslim, but he could just as well mean he was a Nazi.

\paragraph{Lepanto}


Evola discussed the Turkish threat to Europe on one occasion, incorporating both historical details and the philosophy of Nicholas Cusanus. Despite lack of support from the kings of Northern Europe, the Pope was able to rally enough military support to stop the Turks at Lepanto. Interestingly enough, Evola did not praise this. He, along with Cusanus, preferred an ecumenical, or worldwide, unity. Left unsaid, we are left to assume that Evola expects that to happen under a Caliphate.

That will shock the European new right, who for the most part oppose Muslims in Europe. They will never give Pope Pius II credit for his efforts, not will they consider their present plight the result of divine retribution.

Gelli makes much of an alleged prediction by Evola that Islam would someday dominate Europe. If so, he is merely parroting Guenon, which we pointed out already\footnote{\url{https://www.gornahoor.net/?p=35}}. A spiritual vacuum in the West will make that likely, if not inevitable. Paganism has no chance of filling that vacuum, and the traditional spiritual authority of Europe is rejected by Evola; we can see Gelli's point.

\paragraph{Pope Julius}


Evola revealed a dream he had of becoming pope. He replaced the statues of the saints in the Vatican with those of the Roman deities. However, the saints rebelled and restored their statues. In frustration, Pope Julius conquered Moscow, where his troops slaughtered the populace indiscriminately. The Pope approved since he does not love his enemies. Perhaps out of boredom, he asked them to stop. When they refused he returned to Rome, he found it full of Turks who worshipped him as the Mahdi. The new pope balked when they asked him to be circumcised. Evola then woke up.

Readers may take that as a hidden message that paganism is a dead end path. Gelli never took Evola's paganism seriously, regarding it as a means to shock the complacent. Nevertheless, many today take Evola's paganism, racism, and anti-semitism quite seriously, so if that was not his intent, he should have been more clear.

Pagans don't shock anyone anymore. Perhaps had he converted to Traditional Catholicism (as did many of his followers in the 1980s), he would have shocked more people. At least, he would have kept his precious foreskin.

\paragraph{Love and Sex}


Evola curiously asserted that initiates cannot marry and have families. This is contrary to the entire Roman Tradition, which was dominated by fathers, who were the priests of their own household hearths. Instead, he proposes the life of a single man, but not without sexual relations. However, this is possible only in the modern and decadent world, not in the world of Tradition where women were more protected; certainly not in the world of Islam.

In talking of Eros, he points out that lovers are not thinking of children in the heat of passion. Too bad Fr. Gelli did not point out Dante's part of Hell for lovers. Keep reminding me, as I will do that soon.

It is a sign of intelligence to see the connection between sex and reproduction. Obviously, we now have superior technology, in contraception and safe abortions, to block that link. Evola fully supports that. He did not live to see the effects of population depletion now afflicting Europe, the effects of which won't be clear for another generation. However, the Montini pope did see that, but Evola never forgave him for a negative comment made 40 years before.

\paragraph{Loyalty}


Despite the talk of loyalty and keeping one's word, Evola seems loyal to no one or anything. He blames the Italian people for losing the War, calling them shit, and unworthy of their great Fascist state. However, a popular army is not composed of the warrior caste, and modern war is not won by valour alone, but rather by amassing superior and overwhelming resources. The Fascist conception is that the leader creates the people. So, it was the state that was shit, sending Italy into a war it was not materially prepared for.

He was not at all loyal to the spiritual authority of Italy. An anti-nationalist, he at one point praised the spiritual unity of Europe in the Middle Ages. Yet, simultaneously, he rejected that very spirit. That is not the kind of spiritual contradiction the Nicholas Cusanus was describing, but comes closer to incoherence.

\paragraph{Conclusion}


Shortly before Evola's death, Gelli emigrated to the UK, a nation and people that Evola had little respect for. That is where the story ends.

Gelli's recollections are of this nature. There are many interesting tidbits, but nothing in much depth. Gelli assumes Evola took a special liking to him as a ``son''; that makes little sense for reasons I'll explain another time. What seems more likely is that Gelli stood out from the rest of the rightwing pack, so Evola enjoyed him as a conversation partner. Gelli has the intelligence to keep up with Evola's many interests and he writes in an engaging fashion and is a good teller of tales. For those interested in the man behind the ideas, The Sufi of Rome is a satisfactory way to spend a summer afternoon in the mountains with your Kindle.

But, it seems Fr. Gelli has a further purpose in mind based on his rejection of Evola's alleged racism and paganism. These, from his perspective, are merely Evola's way of blame, hiding his true spiritual perspective from the world. Evola would then be the second Rumi (which means ``Roman'' in Arabic), working, perhaps on the supernatural plane, to bring Islam to the world. Or else he may just be a ``beautiful loser'', destined to be misunderstood and ignored. I don't think Fr. Gelli has cleared up any misunderstandings.

Readers will have to decide about Evola, but perhaps that is really Fr. Gelli's projection. He told me he owes much to the spiritual perspective of \textbf{Henri Corbin}. Corbin was exoterically a Protestant, with a degree in Catholic Thomist philosophy. However, Corbin rejects the Incarnation. Corbin devoted much of his life to the study of Shi'ite and Sufi spirituality, of the sort explicitly preferred by Evola.

Corbin created a complex system relating the traditions of pre-Islamic Iran, Shi'ism, Hermetism, and Christianity. He also described a spiritual geography, where the East is not where you think, and the North, what we call Hyperborea, where the Midnight Sun (the Black Sun?) is the true light. Corbin also relates the Iranian cosmology to that of Nordic paganism. The Walkyries are equivalent to the Fravartis. This is just the surface of much deeper correspondences, which are too complex to discuss here.

Is this something known by the Sufi of Rome? Fr. Gelli would have us think so.


\flrightit{Posted on 2012-06-10 by Cologero}

\begin{center}* * *\end{center}

\begin{footnotesize}
\begin{sffamily}
\texttt{Avery on 2012-06-10 at 20:42 said:}

I think Evola's dream about becoming Pope just eliminated whatever remaining romance Classical neopaganism had for me.

\hfill

\texttt{Jason-Adam on 2012-06-11 at 03:49 said:}

I dont think Gelli is right in writing that Evola was against the defence of Europe from the Turks.

In Revolt against the Modern World, he wrote that it was a sign of decline that Christendom was unable to mount a crusade to recover Constantinople after it fell in 1453, also in the same book he castigated Francois I, the king of France, for making an alliance with the Turks and betraying the unity of Europe-Christendom.

In other writings, Evola made it clear he opposed conversion to Islam because he saw Islam as being un-European. It is important to note that while Evola did believe in the unity of Traditions, he did see the outward exoteric form (the culture) as not being transferable. It was not enough for him to be Traditional, he also had to be European.

If anyone is interested, I can post some links to articles by Evola that back up my points.

\hfill

\texttt{apeiron on 2012-06-11 at 22:50 said:}

Since I haven't purchased a kindle and the author appears to have side-stepped a printed version(un-peered?), we'll just have to rely on second or third hand even information from the available reviews and comments. ``The way of blame'' is a fitting eulogy for an apprentice trying to rationalize his masters life and actions? In this case just live as you wish and someone will fill in the gaps of the misunderstood hero when you're dead.

Anyways to what extent can these claims be verified, if at all true:

1)''Evola showed Gelli a letter from an Argentine woman who claimed she gave birth to his son.'' Has Evola a son/s (if the promiscuity is to be believed) somewhere down in Argentina we never knew about?

2)'' ... Gelli often attended their rallies. He even became a Maoist in 1969, so it seems several years of conversation with Evola did not leave a lasting impression on him.''

This is most surprising but not too far-fetched, given Freda's position and Yockey's post-war ``solidarity with the third world'' anti-Anglo/American forces jibe. Evola had referred to those who had ``regressed in character'' later (not because of) despite reading his works. Perhaps Gelli and others fit this profile.

3)Omar Amin von Leers : Could easily be ``on our side'', the semantics are of little importance either way.

4)Lepanto: Jules knew the crescent and the scimitar as lunar symbols of the reflected silver Semite. The solar cross was and is ``our symbol''.

5)Pope Julius: What truth is there to this alleged dream or musing?

6)Love and Sex: ``Evola curiously asserted that initiates cannot marry and have families.'' Why did Guénon do Both (marry, have kids and being initiated)?

7)Loyalty: Evola kept a loyalty all his life to Rome(Amor). He also praised those who were willing to fight on in post-war Italy despite all and secondly that it was not Facism that failed but the Italian people who failed Facism regardless of the trappings of modern combat. Perhaps a fair statement.

8)At this point it's too close to call on the veracity of Gelli's ``tales'' and a thorough reading of the book , if ever it be published is in order. However, all else said, Henri Corbin is beginning to sound very interesting.

\hfill

\texttt{Cologero on 2012-06-12 at 18:01 said:}

Apeiron, I think it can be read on any computer, not just a Kindle, although it may be a little pricey for an ebook. I TOFFT.

To be honest, if it concerned Frank Gelli and anyone else, no one would read it. Where he could have offered some good insights he declined. Specifically, he mentioned discussions about the ``left hand path'' and tantric experiences; these, unfortunately, are omitted from the book.

Claudio Mutti, an Italian convert to Islam, who knew Evola and Omar Amin von Leers wrote an thorough analysis of all of Evola's writings on Islam in \textit{Islam in the eyes of Julius Evola}\footnote{\url{http://www.claudiomutti.com/index.php?url=6&imag=1&id\_news=130}}. Clearly Evola had a high regard for Islam, although it appears doubtful he ``converted'' in any normal sense of the world.

\hfill

\texttt{AghorNath on 2012-06-13 at 02:19 said:}

``he saw Islam as being un-European. It is important to note that while Evola did believe in the unity of Traditions, he did see the outward exoteric form (the culture) as not being transferable. It was not enough for him to be Traditional, he also had to be European.''

On the topic of transference it should be noted that Islamic civilization, wherever it's seeds were sown (by the very fact of its inherent universality and primordiality) always allowed the soil of the native lands which adopted an Islamic identity to retain its traditional weltanschauung, ethnic identity, customs(as long as it did not contradict `tawhid' i.e. Divine Unity) and other localized flavors. The cultural manifestation of Islamic culture adapts where it's winds blow. The various cultural `zoning' of Islamic migration which include, Arabic, Persian, Turkic, African(both in its black and arab/berber manifestations), Malay, Indian, Chinese, Slavic and a few others, only shows that its cultures and arts took on distinct localized character of it's Islamic inheritance. I find it hard to believe that Baron Evola found the idea that European identity would be usurped by any primordial traditional social changes. If perhaps Evola did utter that Islam was ``un-European''(keeping in mind that if he did, he could have said it in his earlier days, and changed his mind as he matured), he probably had in mind Sunni Islam, and not Shi'ite Islam, which in fact is esoteric by its very nature, thus granting, or rather, allowing, an organic transferable acceptance of Islam to Europe as a whole.

\hfill

\texttt{Egemen on 2012-06-13 at 04:21 said:}

Your last link; ``Islam in the eyes of Julius Evola'' is fantastic. Thanks for your sharin'.

\hfill

\texttt{Mihai on 2012-06-14 at 06:01 said:}

I think that this whole thing is just a fantastic story, entirely made-up.

\hfill

\texttt{Jason-Adam on 2012-06-18 at 19:08 said:}

In response to AgorNath -- Evola did not have a complete understanding of Islam, there is another article by Claudio Mutti entitled ``Evola and Nasser'' (see \url{http://www.centrostudilaruna.it/evolaenasser.html}) that shows some errors he made such as assuming that the Wahabi sect (an anti-Traditional cult from the XVIIIth century) was somehow orthodox Muslim. Evola also saw Roman Tradition (and the Roman spiritual race it created) as being superior to others and thus he would not accept receiving tradition, as a Roman writing for Romans, from a non-Roman source. Obviously there are many ways the holes in his logic can be attacked. \textit{Evola and Nasser}\footnote{\url{http://www.claudiomutti.com/index.php?url=6&imag=1&id\_news=206}} is also available in English --- admin \hfill

\texttt{Cologero on 2012-06-22 at 22:59 said:}

Yes, Jason-Adam, there are always two points of view on such a matter; e.g., in the Middle Ages, Muslim scholars regarded European races as unfit for civilization. For example, Said of Toledo wrote:

\begin{quotex}\footnotesize

The other peoples (apart from the Chinese and Turks) that do not cultivate the sciences, resemble rather beasts than men; as regards those that live in the lands of the far North, bordering on the uninhabited part of the globe, the prolonged absence of the sun renders the air cold and the atmosphere in which they live less clear; accordingly they are men of a cold temperament and never reach maturity; they are of great stature and of a white colour, with long and lank hair. But they lack all sharpness or wit and penetration of intellect, and among them predominate ignorance and stupidity, mental blindness, and barbarism.
\end{quotex}


\hfill

\texttt{Cologero on 2012-06-29 at 14:18 said:}

When I wrote this review, the book was about 280,000 on Amazon. It rose to 110,000 in the course of a week. Now it is falling like a rock.

\hfill

\texttt{Aghorable on 2012-06-30 at 03:06 said:}

What cause for surprise? Islam in Europe was absorbed most thoroughly by peoples across the Mare Adriaticum from Italia, starting from around 700 A.H. and surviving there until today. Even the Egypt where Guenon and others found solace from the West was fathered by an Albanian, Mehmet Ali Pasha. Perhaps the winds carried some secret songs to Evola from the Dinaric Heights as he stalked the Italian Alps?

We are in agreement with AghorNath's comment above, and add for clarification that the Bosniaks, for instance, are mostly Sunni Muslims, yet this has not prevented a thoroughly organic adaptation of Islam from taking root there. To conceive of Islam in emulation of the Arabic fashion is indeed not mandatory, and has only become prominent with the rise of Wahhabism, and this is to be regretted. The traveler to Bosnia is saddened by the sight of austere white-washed Saudi-sponsored mosques replacing elaborate old Ottoman mosques that were razed in recent conflict.

\hfill

\texttt{Jason-Adam on 2012-07-02 at 21:30 said:}

@Aghorable

Is not ancestor worship though part of Tradition ? One dilemma that will have to be overcome if you are correct is that if Islam is the true religion, then the ancestors of the European aristocracy who fought and died for Catholicism were in error. If shown to be in error, the aristocracy can lose its entire basis for legitimacy.

\hfill

\texttt{Aghorable on 2012-07-03 at 01:51 said:}

Jason-Adam, Islam, understood in the conventional sense as a particular religion, is not `the true religion' any more than Catholicism or any other. The covenant between God and man remains valid in both religions, and their respective faithful do not maintain their belief in vain. So the Catholic aristocracy of the past maintains its right in God.

The very perspective that involves a comparison of these faiths to discern the truth is, in the last analysis, an exoteric one; the highest Truth effaces them both, together with apparent contradictions, and this is the goal of esoterism. Because of the times that we live in, even those students not in a position to aspire to the actual realization of this as a goal are privy to the didactic statement of its truth and consequences. This constitutes a good theoretical preparation for them, but the privilege also entails a certain responsibility; namely, permitting one's consciousness to acknowledge the existence of a superior universalism (not to be confused with its modern parody, `globalism'). Naturally, if one is not up to the task of such a responsibility, he may always retreat to pure and simple worship according to the laws of his faith; no man should be compelled to wrangle with esoterism, nor do we need them doing so.

Ancestor worship is noble, but it does not simply entail unvarying ancestral emulation, because the descendant does not live under the same stars. Nor can it be asked whether a rule is a part of `Tradition' as such, because there are multiple traditional forms*; but we take it that you refer to that continuity of sympathy by blood that lives on in those men who are not deracinated moderns. So if the question is not about discerning a `true' religion and thus determining obligations, but whether a Catholic becoming Muslim by preference thereby dishonors his ancestors, then the answer must be determined according to perspectives and details. In our case, men of European heritage and students of metaphysic, a parallel question asks whether we accept the charge that it was dishonorable for pagans to become Catholic? Are we to unreservedly regret and lament our Catholic heritage on that account? Some do, but we share no such sentiment. In the plenitude of God, there is, so to speak, room for brothers to face off on the battlefield by day, and to drink and dance alongside the fire by night.

The proud race must be able to assimilate anything extrinsic that it chooses to assimilate, thus making it its own in synthesis. And the function of the highest caste, the men of aether, in such a race must be to guide this assimilation with the luminous knowledge that nothing good is really extrinsic, for what can come to us from outside of the Supreme Being?

*If we could recall the title, we would refer readers to Cologero's article on this matter.

\hfill

\texttt{Aghorable on 2012-07-03 at 02:06 said:}

The article mentioned in the above *note is \textit{``The Absurdity of
Traditionalism''}\footnote{\url{http://www.gornahoor.net/?p=728}}

See also \textit{``Guenon and Islam''}\footnote{\url{http://www.gornahoor.net/?p=4412}}

\hfill

\texttt{Saladin on 2012-10-30 at 20:08 said:}

With all due respect Cologero, this article does not convey the irony adequately! And I have a high opinion of your writing in general.

That aside, HOO's mention of crypto-gynaecocratic had me thinking... Could there be some truth to that?!

\hfill

\texttt{Saladin on 2012-10-30 at 20:15 said:}

Jason, anyone familiar with Evola's work clearly knows that he was not against defending Europe from the Turks. Frankly imagining the Maestro as a Sufi makes me laugh! By the way I am a moslem by birth.

\hfill

\texttt{Cologero on 2012-10-30 at 20:58 said:}

With all due respect, Saladin, if I did not convey the irony adequately, it is most likely because I did not see it.

If you believe there is something ``crypto-gynaecocratic'' about Gornahoor, it is up to you to point it out specifically. As Chicalinda recently said, women do not think logically, so they believe it suffices to just put a label on something and call people names, as Hoo did. It is something else to actually build up a case.

Hoo also accused us of fabricating the non-existent meeting between Evola and Guenon. We painstakingly produced correspondence between them over the course of decades. As is obvious, their relationship was cordial, yet strained. We also see from the letters that in the late 40s, Evola was asking for Guenon's picture because \emph{Evola had no idea what Guenon looked like!} Not did he know Guenon's age. So how could anyone believe they had ever met in person?

I'm still waiting for an apology from HOO and I am very patient.

But that is how a man thinks and acts, he builds his case logically and thoroughly, no matter what you may have been told elsewhere.

\hfill

\texttt{Cologero on 2012-10-30 at 21:02 said:}

Saladin, Evola was not the Maestro, Guenon was the Maestro. Evola himself referred to Guenon in those terms.

\hfill

\texttt{Graham on 2012-10-30 at 23:30 said:}

The last paragraph is astute. A penchant for irony often indicates a deeprooted insincerity.

\hfill

\texttt{Saladin on 2012-10-31 at 21:12 said:}

Regarding the ``crypto-Gynaecocratic'' comment: I am new to this site but one of the very first posts I read was about Bulldog (or maybe it was Pitbull). The writer of that post had put this ``artist'' on the pedestal as a ``real'' man worthy of emulation. This and other posts by the same writer was the cause of my raising the G question.

I am with you on the issue of a Guenon-Evola meeting. Although one can never be certain, the logic of your argument against the possibility of such a meeting is convincing enough. By the way, let me thank you for reproducing/translating such valuable information on this Blog as correspondence between Guenon and Evola.

That a man ``builds his case logically and thoroughly'' I agree! A man also is not afraid to apologize, so maybe HOO will do just that.

\hfill

\texttt{Cologero on 2012-11-01 at 10:50 said:}

``Gynaecocracy'', if it means anything, means rule by women; justification for that cannot be found on Gornahoor.

What ``critics'' imply by that, however, is the false distinction between allegedly lunar and solar traditions. Tradition, to be called such, leads to knowledge of the Absolute, or Tao, since we are in a Taoist frame of mind lately. Any system whose ultimate reality is either the yin or the yang (as the terms lunar and solar would indicate), rather than the Tao, is incomplete, hence cannot be a Tradition in any true sense.

Would anyone seriously consider Dante, Hafiz, Rumi, Ibn Al'Arabi to be crypto-gynaecocratic because of their acknowledgment of the feminine?

I think Chicalinda mentioned pitbull as a man with charisma able to influence women; that is simply a fact. What I would take from it is that a ``real man'' will also have that charisma, even if there are few real men around today. A ``real man'' with such charisma will inspire entire nations to follow God's ways. I believe, Saladin, you may be able to think of some examples.

\hfill

\texttt{Frank Gelli on 2016-12-12 at 11:37 said:}

Haha! Zarathustra laughs out loud!

\hfill

\texttt{Cologero on 2016-12-12 at 12:49 said:}

We are all having a good time here, Fr. Gelli!

\hfill

\end{sffamily}
\end{footnotesize}

